% Brian Ebiyau — CV (industry)
% Build:  cd cv && pdflatex -jobname=Brian_Ebiyau_CV cv.tex
%         then copy Brian_Ebiyau_CV.pdf to ../site/public/
\documentclass[10pt,letterpaper]{article}

\usepackage[hmargin=0.75in, vmargin=0.7in]{geometry}
\usepackage[T1]{fontenc}
\usepackage{microtype}
\usepackage{titlesec}
\usepackage{needspace}
\usepackage{etoolbox}
% Never strand a section heading at the bottom of a page
\pretocmd{\section}{\needspace{6\baselineskip}}{}{}
\usepackage{enumitem}
\usepackage[dvipsnames]{xcolor}
\usepackage[colorlinks, urlcolor=MidnightBlue, linkcolor=black]{hyperref}
\hypersetup{
  pdftitle={Brian Ebiyau -- CV},
  pdfauthor={Brian Ebiyau},
  pdfkeywords={mmWave radar, machine learning, computer vision, PyTorch, quantization, LLM fine-tuning},
}

\pagestyle{empty}
\setlength{\parindent}{0pt}
\linespread{1.08}
\hyphenation{caregiver response}

% Section headings: bold capitals with a rule underneath
\titleformat{\section}{\large\bfseries}{}{0em}{\MakeUppercase}[\vspace{-0.6em}\rule{\linewidth}{0.6pt}]
\titlespacing*{\section}{0pt}{1.5em}{0.7em}

\setlist[itemize]{leftmargin=1.5em, topsep=0.3em, itemsep=0.3em, parsep=0pt}

% \project{Title}{Date}: a one-line heading, followed by bullets.
% Plain lines (no tables) so applicant-tracking systems read dates next to their entries.
\newcommand{\project}[2]{\needspace{4\baselineskip}\textbf{#1}\hfill #2\par}

% \pub{Title}{Link}{Venue}{Authors}
\newcommand{\pub}[4]{\textbf{\href{#2}{\color{black}#1}}, #3\\ #4\par\vspace{0.6em}}

\newcommand{\me}{\textbf{Brian Ebiyau}}

\begin{document}

% ---------------------------------------------------------------- Header
\begin{center}
  {\LARGE\bfseries Brian Ebiyau}\\[0.35em]
  \small
  Pittsburgh, PA \enspace|\enspace
  \href{mailto:bebiyau@andrew.cmu.edu}{bebiyau@andrew.cmu.edu} \enspace|\enspace
  \href{https://www.linkedin.com/in/brian-ebiyau/}{linkedin.com/in/brian-ebiyau} \enspace|\enspace
  \href{https://github.com/ebyau}{github.com/ebyau} \enspace|\enspace
  \href{https://scholar.google.com/citations?user=pZlE07AAAAAJ&hl=en}{Google Scholar}
\end{center}

% ---------------------------------------------------------------- Research focus
\section{Research Focus}
Building toward foundation perception models that sense and understand the world through mmWave radar, to establish
radar as a scalable, low-cost and robust complement or alternative to cameras and lidar for autonomous systems,
robotics and smart environments. Seeking perception internships.

% ---------------------------------------------------------------- Education
\section{Education}
\textbf{Carnegie Mellon University}, Ph.D., Electrical and Computer Engineering \hfill Aug. 2026 -- Expected 2031\par
Advisors: Anthony Rowe, Jaime Cardoso\par\vspace{0.45em}
\textbf{Carnegie Mellon University}, M.S., Engineering Artificial Intelligence \hfill 2023 -- 2025\par\vspace{0.45em}
\textbf{Makerere University}, B.S., Electrical and Computer Engineering \hfill 2017 -- 2022\par

% ---------------------------------------------------------------- Skills
\section{Skills}
\textbf{Radar Perception \& Sensing:} 4D mmWave radar, radar signal processing, TDM/DDM MIMO systems\\[0.3em]
\textbf{Machine Learning \& Foundation Models:} Deep learning, Transformers, foundation models, LoRA fine-tuning\\[0.3em]
\textbf{Computer Vision \& Multimodal Perception:} Object detection and tracking, privacy-preserving vision systems\\[0.3em]
\textbf{Signal Processing:} Range-Doppler-Angle processing, FFT-based analysis, CFAR detection, point cloud generation\\[0.3em]
\textbf{Programming \& AI Frameworks:} Python, PyTorch

% ---------------------------------------------------------------- Research experience
\section{Research Experience}
\project{RadarML Group, WiseLab, Carnegie Mellon University \textmd{--- \textit{Graduate Researcher}}}{Aug. 2026 -- Present}
\begin{itemize}
  \item Explored range-dependent tokenization of raw 4D radar for 3D occupancy prediction with WiseLab's
        \href{https://wiselabcmu.github.io/grt/}{Generalizable Radar Transformer (GRT)}, allocating finer tokens to far
        range; benchmarked against uniform patching on I/Q-1M.
  \item Implementing Doppler-division multiplexing (DDM) MIMO in place of time-division
        multiplexing (TDM) on the TI AWR1843AOP EVM, enabling simultaneous transmission from all
        TX antennas.
\end{itemize}
\vspace{0.6em}
\project{Mobile Edge AI Lab, Carnegie Mellon University Africa \textmd{--- \textit{Researcher}}}{2025 -- 2026}
\begin{itemize}
  \item Led the first large-scale study of driver responses to motorcycles in African urban traffic (36,658
        encounters from 568 dashcam videos in Kigali; first author,
        \href{https://ieeexplore.ieee.org/abstract/document/11376493/}{IEEE ICVES 2025}), identifying three behavioral
        strategies and reaction-time baselines for ADAS/AV calibration.
\end{itemize}

% ---------------------------------------------------------------- Projects
\section{Projects}
\project{Stereo Face Anonymization for Multi-Camera Video}{2026}
\begin{itemize}
  \item Anonymized 100+ hours of multi-camera video for open research: localized faces from precomputed pose
        keypoints in the left camera, projected them to the right via stereo disparity, and re-encoded videos with
        metadata intact.
\end{itemize}
\vspace{0.6em}
\project{Domain Adaptation of Small LLMs for Telecom QA}{2024}
\begin{itemize}
  \item Adapted Phi-2 for QA on 3GPP telecom standards using LoRA fine-tuning and context-response conditioning,
        winning \href{https://zindi.world/competitions/specializing-large-language-models-for-telecom-networks/leaderboard}{\textbf{1st place among 50+ teams}} with 81.9\% accuracy; code and models are
        \href{https://github.com/Alexgichamba/itu_qna_challenge}{open source}.
\end{itemize}

% ---------------------------------------------------------------- Work experience
\section{Work Experience}
\project{Kin-Keepers, USA \textmd{--- \textit{Intern}}}{Jul. 2021 -- Jul. 2022}
\begin{itemize}
  \item Developed Alzheimer's screening models from patient voice samples using TF-IDF features with
        MLP and logistic regression, reaching $\sim$80\% accuracy on DementiaBank while
        screening 5$\times$ faster than pen-and-paper tests.
\end{itemize}

% ---------------------------------------------------------------- Publications
\section{Publications}
% ICVES paper omitted from Publications for now (linked from the Mobile Edge AI Lab bullet).
\pub{Eyes and Ears: Automated Annotation of Audio Data Using Computer Vision}
    {https://ieeexplore.ieee.org/abstract/document/11084628/}
    {IEEE ICIP 2025}
    {Galane Basha Namomsa, Alex Gichamba, \me, João Barros}
\pub{ColBERT Retrieval and Ensemble Response Scoring for Language Model Question Answering}
    {https://doi.org/10.1109/GCWkshp64532.2024.11100297}
    {IEEE GLOBECOM Workshops 2024}
    {Alex Gichamba, Tewodros Kederalah Idris, \me, Eric Nyberg, Teruko Mitamura}
\vspace{-0.6em}

% Honors & Awards omitted for now (see archive/cv-industry-v11.tex).

% ---------------------------------------------------------------- Teaching
\section{Teaching}
\textbf{Teaching Assistant}, CMU 11-785 Introduction to Deep Learning \hfill Spring 2025\par\vspace{0.3em}
\textbf{Guest Lecturer}, CMU 11-611 Natural Language Processing \hfill Fall 2024\par

% ---------------------------------------------------------------- Service
\section{Service}
\textbf{Reviewer}, IEEE International Conference on Intelligent Transportation Systems (ITSC) \hfill 2026\par\vspace{0.3em}
\textbf{Member}, Ph.D. in Education Student Organization (PESO), CMU ECE \hfill 2026\par

% Selected coursework (removed for industry; restore if wanted):
% Computer Vision, Deep Learning, Natural Language Processing, Radar Machine Learning

\end{document}
